% ============================================================
\chapter{M\'ethodes de Points Int\'erieurs}
\label{ch:points_interieurs}
% ============================================================

\section{Introduction}

En 1984, Narendra Karmarkar, ing\'enieur chez AT\&T Bell Labs, publie un algorithme qui secoue le monde de l'optimisation~: une m\'ethode de programmation lin\'eaire en temps polynomial qui, contrairement au simplexe de Dantzig, ne parcourt pas les sommets du poly\`edre mais traverse son \emph{int\'erieur}. La communaut\'e est \'electris\'ee --- et divis\'ee. Certains voient une r\'evolution, d'autres un coup m\'ediatique (le simplexe, bien que de complexit\'e exponentielle dans le pire cas, est rapide en pratique). Mais l'histoire tranchera~: les m\'ethodes de points int\'erieurs, raffin\'ees par Nesterov et Nemirovski dans les ann\'ees 1990, s'av\`erent d\'ecisives pour l'optimisation semi-d\'efinie, l'optimisation conique, et les probl\`emes de grande taille. L'id\'ee est \'el\'egante~: remplacer les contraintes d'in\'egalit\'e par une barri\`ere logarithmique, puis suivre le \emph{chemin central} vers l'optimum.

\begin{intuition}
Imaginez que vous \^etes dans une pi\`ece (le poly\`edre des contraintes) et
que vous cherchez le point le plus bas (l'optimum). Le simplexe longe les murs
et les coins. Les m\'ethodes de points int\'erieurs, elles, marchent au milieu
de la pi\`ece en se rapprochant progressivement de la solution optimale, qui
se trouve souvent pr\`es d'un mur. Une barri\`ere invisible vous emp\^eche
de toucher les murs, mais cette barri\`ere s'affaiblit au fil des it\'erations.
\end{intuition}

\section{M\'ethode de barri\`ere}

\subsection{Formulation}

\begin{definition}[Probl\`eme avec in\'egalit\'es]
On consid\`ere le probl\`eme convexe~:
\[
  \min_{x \in \R^n} f_0(x) \quad \text{s.c.} \quad
  f_i(x) \leq 0, \; i = 1, \ldots, m, \qquad
  Ax = b,
\]
o\`u $f_0, f_1, \ldots, f_m$ sont convexes et deux fois diff\'erentiables.
\end{definition}

\begin{definition}[Fonction barri\`ere logarithmique]
\index{barri\`ere logarithmique}
La \textbf{fonction barri\`ere logarithmique} est~:
\[
  \phi(x) = -\sum_{i=1}^m \ln(-f_i(x)),
\]
d\'efinie sur l'int\'erieur strict du domaine r\'ealisable
$\{x : f_i(x) < 0, \; \forall i\}$.
\end{definition}

\begin{definition}[Probl\`eme barri\`ere]
\index{m\'ethode de barri\`ere}
Le \textbf{probl\`eme barri\`ere} pour un param\`etre $t > 0$ est~:
\[
  \min_{x} \; t \, f_0(x) + \phi(x) \quad \text{s.c.} \quad Ax = b.
\]
On note $x^*(t)$ la solution, appel\'ee \textbf{point central} pour le
param\`etre $t$.
\end{definition}

\subsection{Chemin central}

\begin{definition}[Chemin central]
\index{chemin central}
Le \textbf{chemin central} est l'ensemble des points
$\{x^*(t) : t > 0\}$. Quand $t \to +\infty$, $x^*(t) \to x^*$, la
solution du probl\`eme original.
\end{definition}

\begin{theoreme}[Borne de sous-optimalit\'e]
Pour tout $t > 0$, le point central $x^*(t)$ v\'erifie~:
\[
  f_0(x^*(t)) - p^* \leq \frac{m}{t},
\]
o\`u $p^*$ est la valeur optimale et $m$ le nombre de contraintes
d'in\'egalit\'e.
\end{theoreme}

\begin{proof}
Les conditions KKT du probl\`eme barri\`ere montrent que $x^*(t)$ est
r\'ealisable pour le probl\`eme original, avec des multiplicateurs
$\lambda_i^*(t) = -1/(t f_i(x^*(t)))$. Par la dualit\'e faible~:
\[
  f_0(x^*(t)) - p^* \leq \sum_{i=1}^m \lambda_i^*(t) \cdot (-f_i(x^*(t)))
  = \sum_{i=1}^m \frac{1}{t} = \frac{m}{t}.
\]
\end{proof}

\begin{algorithme}[title=\textbf{M\'ethode de barri\`ere}]
\begin{enumerate}
  \item Choisir $x_0$ strictement r\'ealisable, $t_0 > 0$,
        facteur $\mu > 1$, tol\'erance $\varepsilon > 0$.
  \item Pour $k = 0, 1, 2, \ldots$~:
    \begin{enumerate}
      \item \textbf{Centrage}~: r\'esoudre (par Newton)
            $\min_x t_k f_0(x) + \phi(x)$ s.c. $Ax = b$,
            en partant de $x_k$, pour obtenir $x_{k+1} = x^*(t_k)$.
      \item \textbf{Test d'arr\^et}~: si $m / t_k < \varepsilon$,
            \textbf{stop}.
      \item \textbf{Mise \`a jour}~: $t_{k+1} = \mu \cdot t_k$.
    \end{enumerate}
\end{enumerate}
\end{algorithme}

\begin{formules}[title=\textbf{Complexit\'e de la m\'ethode de barri\`ere}]
\begin{itemize}
  \item Nombre d'it\'erations ext\'erieures (pas de barri\`ere)~:
        $\lceil \frac{\ln(m / (t_0 \varepsilon))}{\ln \mu} \rceil$.
  \item Choix typique~: $\mu = 10$ \`a $20$, ce qui donne peu d'it\'erations
        ext\'erieures.
  \item Chaque \'etape de centrage n\'ecessite quelques pas de Newton (6--40
        en pratique).
\end{itemize}
\end{formules}

\section{M\'ethodes \`a pas court et \`a pas long}

\begin{definition}[M\'ethode \`a pas court]
\index{m\'ethode \`a pas court}
La m\'ethode \`a \textbf{pas court} choisit $\mu$ proche de 1
(par exemple $\mu = 1 + 1/\sqrt{m}$) et effectue un seul pas de Newton
par \'etape de centrage. Le nombre d'it\'erations est
$\mathcal{O}(\sqrt{m} \ln(m/\varepsilon))$.
\end{definition}

\begin{definition}[M\'ethode \`a pas long]
\index{m\'ethode \`a pas long}
La m\'ethode \`a \textbf{pas long} choisit $\mu$ grand et effectue
plusieurs pas de Newton pour recentrer. Le nombre total de pas de Newton est
aussi $\mathcal{O}(\sqrt{m} \ln(m/\varepsilon))$ mais avec une constante
plus grande.
\end{definition}

\begin{theoreme}[Complexit\'e polynomiale]
\index{complexit\'e polynomiale}
Les m\'ethodes de points int\'erieurs r\'esolvent un programme lin\'eaire
\`a $n$ variables et $m$ contraintes en $\mathcal{O}(\sqrt{m} \ln(1/\varepsilon))$
pas de Newton, chacun co\^utant $\mathcal{O}(n^3)$ (ou moins avec
l'exploitation de la structure creuse). La complexit\'e totale est donc
polynomiale.
\end{theoreme}

\section{M\'ethode primale-duale}

\begin{definition}[Syst\`eme KKT modifi\'e]
La m\'ethode \textbf{primale-duale} r\'esout simultan\'ement les conditions
KKT perturb\'ees~:
\[
  \begin{cases}
    \nabla f_0(x) + \sum_i \lambda_i \nabla f_i(x) + A^\top \nu = 0, \\
    -\lambda_i f_i(x) = 1/t, \quad i = 1, \ldots, m, \\
    Ax = b.
  \end{cases}
\]
On applique la m\'ethode de Newton \`a ce syst\`eme.
\end{definition}

\begin{remarque}
La m\'ethode primale-duale est plus efficace en pratique que la m\'ethode de
barri\`ere pure car elle ne n\'ecessite pas de r\'esoudre exactement le
probl\`eme de centrage \`a chaque \'etape. Les solveurs modernes (MOSEK,
CVXOPT) utilisent cette approche.
\end{remarque}

\begin{theoreme}[Convergence primale-duale]
La m\'ethode primale-duale atteint une pr\'ecision $\varepsilon$ en
$\mathcal{O}(\sqrt{m} \ln(1/\varepsilon))$ it\'erations de Newton, sous
des hypoth\`eses standards de r\'egularit\'e.
\end{theoreme}

\section{Application \`a la programmation lin\'eaire}

\begin{exemple}[PL sous forme standard]
Pour le programme lin\'eaire $\min c^\top x$ s.c. $Ax = b$, $x \geq 0$,
la fonction barri\`ere est~:
\[
  \phi(x) = -\sum_{j=1}^n \ln(x_j).
\]
Le syst\`eme de Newton pour le centrage s'\'ecrit~:
\[
  \begin{pmatrix}
    X^{-2} & A^\top \\
    A & 0
  \end{pmatrix}
  \begin{pmatrix} \Delta x \\ \Delta \nu \end{pmatrix}
  = -\begin{pmatrix} t c - X^{-1} e \\ 0 \end{pmatrix},
\]
o\`u $X = \mathrm{diag}(x)$ et $e = (1, \ldots, 1)^\top$.
\end{exemple}

\section{Exercices}

\begin{exercice}[$\star$ -- Barri\`ere logarithmique]
Pour le PL $\min -x_1 - x_2$ s.c. $x_1 + x_2 \leq 1$, $x_1, x_2 \geq 0$,
\'ecrire la fonction barri\`ere et calculer $x^*(t)$ pour $t = 1, 10, 100$.
V\'erifier que $x^*(t) \to x^*$.
\end{exercice}

\begin{exercice}[$\star\star$ -- Chemin central]
Tracer le chemin central pour le probl\`eme $\min x_1$ s.c.
$x_1^2 + x_2^2 \leq 1$ pour $t \in \{0.1, 1, 10, 100\}$. Montrer qu'il
converge vers $(-1, 0)$.
\end{exercice}

\begin{exercice}[$\star\star$ -- Pas de Newton]
D\'eriver le syst\`eme de Newton pour la m\'ethode de barri\`ere appliqu\'ee
au probl\`eme quadratique contraint~:
$\min \frac{1}{2} x^\top Q x + c^\top x$ s.c. $Ax \leq b$.
\end{exercice}

\begin{exercice}[$\star\star\star$ -- Borne de sous-optimalit\'e]
D\'emontrer en d\'etail que $f_0(x^*(t)) - p^* \leq m/t$. En d\'eduire le
nombre d'it\'erations ext\'erieures n\'ecessaires pour atteindre une
pr\'ecision $\varepsilon$ avec facteur $\mu$.
\end{exercice}

\begin{exercice}[$\star\star\star$ -- Primale-duale]
Impl\'ementer la m\'ethode primale-duale de points int\'erieurs pour le PL
sous forme standard. Tester sur un probl\`eme \`a 50 variables et 20
contraintes. Comparer le nombre d'it\'erations avec la m\'ethode de
barri\`ere classique.
\end{exercice}
